\chapter{The Asset Manager and the Fund}\label{ch:fm:the-asset-manager-and-the-fund}

On 3 June 2019 a UK equity income fund suspended dealing. It had been worth more than £10.1 billion two years earlier and was worth £3.6 billion
in the run-up to the suspension; its investors had kept asking for their money back, and its managers had kept selling what was easy to sell.
The UK regulator later found that at the suspension only 8\% of the fund's investments could be sold within seven days, while its investors
were entitled to their money within four. A fund's liquidity terms are a promise about its assets, written in the documents of a vehicle; this
chapter reads the vehicle, its fees and its promises as the manager must run them.

\section{Structures: manager, fund and vehicle}

An asset manager (One Quant Book 1, chapter 1) does not own the money it manages. It runs a fund, a separate legal vehicle owned by its investors,
under a mandate (Book 1, chapter 3), and is paid fees out of it. The separation is what protects investors if the manager fails, and what lets
one strategy be sold in several wrappers.

\begin{definition}[General partner, limited partner]\label{def:fm:the-asset-manager-and-the-fund:gp}
Many private funds are limited partnerships. The \emph{general partner}\index{general partner} manages the partnership, is liable for its debts
and usually delegates investment decisions to the manager, an affiliate. The \emph{limited partners}\index{limited partner} are the investors:
they contribute capital, share in its gains and losses, and are liable only up to what they contributed.
\end{definition}

\begin{definition}[Master--feeder structure]\label{def:fm:the-asset-manager-and-the-fund:master}
A \emph{master--feeder structure}\index{master--feeder structure} is a fund in which investors buy into one or more feeder funds, each built
for a group of investors (onshore taxable investors, offshore or tax-exempt ones), and the feeders invest all their assets in a single master
fund, which holds the portfolio.
\end{definition}

\begin{definition}[UCITS, separately managed account]\label{def:fm:the-asset-manager-and-the-fund:ucits}
A \emph{UCITS}\index{UCITS} is a fund authorised under the European Union's directive on undertakings for collective investment in
transferable securities: it may be sold to the public across the Union, and in exchange it must invest within eligible assets and
diversification limits and redeem its units at any investor's request, suspending only in exceptional cases. A \emph{separately managed
account}\index{separately managed account} is a portfolio held in one investor's own name and managed under an investment management agreement,
outside any pooled fund.
\end{definition}

\begin{omfigure}
\begin{tikzpicture}[font=\footnotesize,>=stealth,box/.style={draw,rounded corners,minimum height=0.62cm,align=center,fill=omDef!10},
  svc/.style={draw,rounded corners,minimum height=0.62cm,align=center,fill=gray!12}]
\node[box,minimum width=2.8cm] (on) at (-2.2,3.4) {onshore feeder\\{\scriptsize taxable investors}};
\node[box,minimum width=2.8cm] (off) at (2.2,3.4) {offshore feeder\\{\scriptsize non-resident, tax-exempt}};
\node[box,minimum width=3.0cm,fill=omProp!12] (m) at (0,1.8) {master fund\\{\scriptsize holds the portfolio}};
\node[box,minimum width=2.4cm,fill=omThm!10] (mgr) at (-4.6,1.8) {manager\\{\scriptsize investment decisions}};
\node[svc,minimum width=2.2cm] (adm) at (-3.3,0) {administrator\\{\scriptsize NAV, register}};
\node[svc,minimum width=2.2cm] (cus) at (-0.9,0) {custodian\\{\scriptsize safekeeping}};
\node[svc,minimum width=2.2cm] (pb) at (1.5,0) {prime broker\\{\scriptsize financing}};
\node[svc,minimum width=2.2cm] (aud) at (3.9,0) {auditor\\{\scriptsize annual accounts}};
\draw[->] (on) -- (m); \draw[->] (off) -- (m);
\draw[->,omThm] (mgr) -- node[above,font=\scriptsize] {mandate} (m);
\foreach \s in {adm,cus,pb,aud} {\draw[<->,gray] (m.south) -- (\s.north);}
\end{tikzpicture}

\omcaption{A master--feeder hedge fund. Investors buy into the feeder built for them; the feeders invest in the master, which holds the
portfolio. The manager decides; the administrator computes the net asset value and keeps the register, the \omterm{def:fm:the-asset-manager-and-the-fund:admin}{custodian} holds the assets, the prime
broker finances the positions, and the auditor signs the accounts.}\label{fig:fm:the-asset-manager-and-the-fund:structure}
\end{omfigure}

\section{Service providers: administrator, custodian, auditor, prime broker}

\begin{definition}[Fund administrator, custodian]\label{def:fm:the-asset-manager-and-the-fund:admin}
A \emph{fund administrator}\index{fund administrator} is the firm, independent of the manager, that computes a fund's net asset value, keeps
its register of investors, processes subscriptions and redemptions and calculates each investor's fees. A \emph{custodian}\index{custodian} is
the institution that holds a fund's assets in safekeeping, in accounts segregated from its own and from the manager's.
\end{definition}

The service providers are the fund's controls. An administrator that prices the portfolio independently stops a manager marking its own book; a
\omterm{def:fm:the-asset-manager-and-the-fund:admin}{custodian} that holds the assets stops them being used for anything but the fund; an auditor tests both once a year. A prime broker (Book 1,
chapter 6) finances the positions and often holds the assets it lends against, which is why a hedge fund's documents say what may be
rehypothecated (chapter 18). A manager chooses these firms, but investors' due diligence (chapter 25) checks each of them.

\section{Fees in practice: classes, hurdles, crystallisation and equalisation}

A fund's fees are the management and performance fees of Book 1, chapter 3, with a high-water mark; in practice they are set per \omterm{def:fm:the-asset-manager-and-the-fund:class}{share class},
and three mechanisms decide who pays what.

\begin{definition}[Share class, fee hurdle, crystallisation]\label{def:fm:the-asset-manager-and-the-fund:class}
A \emph{share class}\index{share class} is a category of a fund's shares with its own terms: currency, fees, liquidity, minimum investment. A
\emph{fee hurdle}\index{fee hurdle} is a return, fixed or tied to a benchmark or a cash rate, that the fund must exceed before a performance fee
is due. \emph{Crystallisation}\index{crystallisation} is the moment at which an accrued performance fee becomes payable to the manager and can no
longer be reversed by later losses; the fund's documents set its frequency.
\end{definition}

The more often fees crystallise, the more the manager earns on the same returns: a gain crystallised in March is not given back when April loses.
On five years of monthly returns with a mean of 8\% a year and a volatility of 12\%, a 20\% performance fee crystallised every month takes 11.54\%
of the starting capital on average over 1\,000 paths, against 11.26\% quarterly and 10.78\% annually.

\begin{definition}[Equalisation]\label{def:fm:the-asset-manager-and-the-fund:eq}
\emph{Equalisation}\index{equalisation} is the set of adjustments, or the practice of issuing a separate series of shares for each subscription
date, by which each investor in a pooled fund pays performance fees on the gains of its own investment, measured from its own entry price.
\end{definition}

\begin{example}[The free ride]\label{ex:fm:the-asset-manager-and-the-fund:free}
Investor A subscribes 100 at a net asset value of 1.00. The fund falls 10\%; investor B subscribes 100 at 0.90. The fund then rises to 1.05. With
one pooled share price and one high-water mark at 1.00, the fund charges 20\% of the 0.05 per share above the mark: A pays 1.00, B pays 1.11. B's
own gain was from 0.90 to 1.05; with series accounting B pays 3.33. Without \omterm{def:fm:the-asset-manager-and-the-fund:eq}{equalisation} B rode free on 2.22 of performance fee that A's loss had
created. The opposite case, an investor who buys above the mark and pays on gains that are not his, is as unfair and as common.
\end{example}

\begin{method}[Checking a fund's fee terms]\label{met:fm:the-asset-manager-and-the-fund:fees}
\begin{enumerate}
\item List every \omterm{def:fm:the-asset-manager-and-the-fund:class}{share class} and its management fee, performance fee, hurdle, high-water mark and \omterm{def:fm:the-asset-manager-and-the-fund:class}{crystallisation} frequency.
\item Check whether performance fees are computed per investor (series or \omterm{def:fm:the-asset-manager-and-the-fund:eq}{equalisation}) or on the pooled share price.
\item Simulate the fee take over paths of the strategy's returns, not on its expected return: \omterm{def:fm:the-asset-manager-and-the-fund:class}{crystallisation} and the mark act on the path.
\item Compare the result with the offering document's worked example; the administrator computes the real fees and should reproduce it.
\end{enumerate}
\end{method}

\section{Liquidity terms and the liquidity mismatch}

\begin{definition}[Redemption notice period, lock-up period]\label{def:fm:the-asset-manager-and-the-fund:notice}
A fund's \emph{redemption notice period}\index{redemption notice period} is the time an investor must give between asking for its money and
the dealing date at which it is redeemed. A \emph{lock-up period}\index{lock-up period} is an initial period after a subscription during which
it cannot be redeemed, or only with a penalty.
\end{definition}

\begin{definition}[Redemption gate, side pocket, swing pricing]\label{def:fm:the-asset-manager-and-the-fund:gate}
A \emph{redemption gate}\index{redemption gate} limits the share of a fund's (or an investor's) assets that can be redeemed on one dealing date;
requests above it are deferred. A \emph{side pocket}\index{side pocket} is a separate account, or \omterm{def:fm:the-asset-manager-and-the-fund:class}{share class}, into which illiquid or
hard-to-value assets are moved, redeemable only when they are sold, so that investors who come and go do not trade them at a guessed price.
\emph{Swing pricing}\index{swing pricing} adjusts the price at which investors subscribe or redeem on days of large net flows, so that the
cost of trading for them falls on them and not on the investors who stay.
\end{definition}

\begin{definition}[Liquidity mismatch]\label{def:fm:the-asset-manager-and-the-fund:mismatch}
A fund has a \emph{liquidity mismatch}\index{liquidity mismatch} when its investors may withdraw a larger share of its value within some period
than it could sell within that period at a normal cost.
\end{definition}

Form PF, the confidential report large US hedge-fund advisers file, gives the aggregate picture: investors in qualifying hedge funds could
withdraw 9.0\% of their capital within seven days in the third quarter of 2025, while the funds reported that they could sell 51.8\% of their
portfolios within seven days (\cref{fig:fm:the-asset-manager-and-the-fund:liquidity}). At every horizon the portfolios were more liquid than the
investors' terms: in aggregate the hedge-fund industry is not mismatched. The fund suspended in 2019 was mismatched by a factor of more than twelve
at seven days: its investors could have all of their money within four days, and it could sell 8\% of its assets within seven.

\begin{omfigure}
\begin{tikzpicture}
\begin{axis}[width=10.5cm,height=5.6cm,xmode=log,log basis x=10,xmin=0.8,xmax=450,ymin=0,ymax=105,
  xtick={1,7,30,90,180,365},xticklabels={1,7,30,90,180,365},xlabel={\small days (log scale)},ylabel={\small cumulative share of NAV (\%)},
  tick label style={font=\footnotesize},grid=major,grid style={gray!20},table/col sep=comma,
  legend style={font=\scriptsize,at={(0.5,-0.3)},anchor=north,/tikz/every even column/.append style={column sep=8pt}},legend columns=2]
\addplot[omThm,thick,mark=*] table[x=days,y=investor]{figdata/desk/04-the-asset-manager-and-the-fund/liquidity.csv};
\addplot[omDef,thick,mark=square*] table[x=days,y=portfolio]{figdata/desk/04-the-asset-manager-and-the-fund/liquidity.csv};
\addplot[gray,thick,dashed,mark=triangle*] table[x=days,y=illiquid]{figdata/desk/04-the-asset-manager-and-the-fund/liquidity.csv};
\addplot[only marks,mark=diamond*,mark size=3pt,omProp] coordinates {(4,100)};
\legend{hedge-fund investors can withdraw,hedge-fund portfolios can be sold,{illiquid fund (8\% in 7 days)},{its investors: all within 4 days}}
\end{axis}
\end{tikzpicture}

\omcaption{Liquidity of the investors' terms against liquidity of the portfolio. US qualifying hedge funds in aggregate, 2025Q3 (Form PF,
SEC Private Funds Statistics, tables 8.22 and 8.23); and an illustrative fund with 8\% of its assets sellable within seven days, the figure the UK
regulator gave for the fund suspended in 2019, whose investors could redeem everything within four days.}\label{fig:fm:the-asset-manager-and-the-fund:liquidity}
\end{omfigure}

\begin{dated}{2026-09}{The public record on fund liquidity}\label{dat:fm:the-asset-manager-and-the-fund:pf}
SEC Private Funds Statistics, third quarter of 2025 (qualifying hedge funds, aggregate NAV \$4\,750 billion): assets that may be suspended \$3\,182
billion and that may be gated \$2\,053 billion; assets gated \$121 billion, side-pocketed \$103 billion and suspended \$25 billion (2.5\%, 2.2\% and
0.5\% of NAV). \omterm{def:fm:the-asset-manager-and-the-fund:ucits}{UCITS} Directive 2009/65/EC, article 84: a \omterm{def:fm:the-asset-manager-and-the-fund:ucits}{UCITS} redeems at any unit-holder's request and may suspend only temporarily, in
exceptional cases and in the unit-holders' interest. US open-end funds, SEC Rule 22e-4: an illiquid investment is one that cannot be sold within
seven calendar days without significantly moving its price, and a fund may not buy one if more than 15\% of its net assets would then be
illiquid; Rule 22c-1(a)(3) permits \omterm{def:fm:the-asset-manager-and-the-fund:gate}{swing pricing} above a threshold of net flows.
\end{dated}

\begin{dated}{2026-09}{The 2019 suspension in the regulator's words}\label{dat:fm:the-asset-manager-and-the-fund:weif}
Financial Conduct Authority, 11 April 2024: the fund's authorised corporate director ``failed to manage the liquidity of the fund'' between
31 July 2018 and the suspension on 3 June 2019; investors in the fund at the suspension receive a share of an up to £230 million redress scheme
approved by the High Court in February 2024. FCA, 5 August 2025: the regulator decided to fine the manager £40 million and its founder
£5\,888\,800 and to ban him from managing retail funds; both referred the decisions to the Upper Tribunal, so the findings are provisional. The
regulator found that from July 2018 the fund ``disproportionately sold more liquid investments'' and bought less liquid ones.
\end{dated}

\begin{proposition}[Liquid first leaves the stayers illiquid]\label{prop:fm:the-asset-manager-and-the-fund:liquidfirst}
Meet a redemption of a share $q$ of NAV by selling the most liquid assets first. If a share $\ell_h$ of NAV can be sold within horizon $h$ and
$q\le\ell_h$, the remaining investors hold a fund of which only $(\ell_h-q)/(1-q)$ can be sold within $h$; if $q\ge\ell_h$, none.
\end{proposition}

\begin{proof}
The sale removes $q$ from the liquid bucket and nothing from the others; the remaining NAV is $1-q$ before costs.
\end{proof}

Selling liquid first minimises the cost of today's redemption and hands the remaining investors a less liquid fund; the next redemption is then
more expensive, and redeeming early becomes worth more than staying. The model of \texttt{firm.fundterms} (\cref{lst:fm:the-asset-manager-and-the-fund:redeem})
meets a 30\% redemption within a seven-day notice period for two ladders, the aggregate hedge-fund portfolio of Form PF and the illiquid fund, with
illustrative selling costs that double for assets sold faster than their horizon:

\begin{center}
\small
\begin{tabular}{@{}llrrr@{}}
\toprule
fund & policy & paid & cost borne by stayers & sellable in 7 days after\\
\midrule
hedge-fund ladder & liquid first & 30\% & 0.04\% & 31.1\%\\
 & pro rata & 30\% & 5.39\% & 51.8\%\\
 & gate at 10\% & 10\% & 0.01\% & 46.4\%\\
 & liquid first, \omterm{def:fm:the-asset-manager-and-the-fund:gate}{swing pricing} & 30\% & 0 & 31.1\%\\
illiquid fund & liquid first & 30\% & 1.69\% & 0\\
 & pro rata & 30\% & 8.81\% & 8.0\%\\
 & gate at 10\% & 10\% & 0.09\% & 0\\
 & liquid first, \omterm{def:fm:the-asset-manager-and-the-fund:gate}{swing pricing} & 30\% & 0 & 0\\
\bottomrule
\end{tabular}
\end{center}

\omcode{code/firm/fundterms/firm_fundterms.py}{85}{112}{Meeting a redemption: liquid first or a slice of everything, a gate, and swing pricing.}\label{lst:fm:the-asset-manager-and-the-fund:redeem}

\begin{omfigure}
\begin{tikzpicture}
\begin{axis}[width=11cm,height=5.4cm,ybar,bar width=10pt,ymin=0,ymax=60,ylabel={\small \%},
  xtick={0,1,2,3},xticklabels={liquid first,pro rata,gate 10\%,{liquid first; swing}},tick label style={font=\footnotesize},
  table/col sep=comma,enlarge x limits=0.15,grid=major,grid style={gray!20},
  legend cell align=left,legend style={font=\scriptsize,at={(0.5,-0.25)},anchor=north,/tikz/every even column/.append style={column sep=8pt}},legend columns=2]
\addplot[fill=omDef!55,draw=omDef] table[x=k,y=hf_liquid]{figdata/desk/04-the-asset-manager-and-the-fund/stress.csv};
\addplot[fill=gray!40,draw=gray] table[x=k,y=ill_liquid]{figdata/desk/04-the-asset-manager-and-the-fund/stress.csv};
\addplot[fill=omThm!55,draw=omThm] table[x=k,y=hf_dilution]{figdata/desk/04-the-asset-manager-and-the-fund/stress.csv};
\addplot[fill=omProp!55,draw=omProp] table[x=k,y=ill_dilution]{figdata/desk/04-the-asset-manager-and-the-fund/stress.csv};
\legend{hedge-fund ladder: sellable in 7 days after,illiquid fund: sellable in 7 days after,hedge-fund ladder: cost to stayers,illiquid fund: cost to stayers}
\end{axis}
\end{tikzpicture}

\omcaption{A 30\% redemption request with seven days' notice, under four policies: the share of the remaining fund sellable within seven days
afterwards, and the selling cost borne by the investors who stay, in per cent of their NAV. Ladders: the Form PF aggregate and the illustrative
illiquid fund; selling costs illustrative. Data: \texttt{fm\_fund.stress\_table}.}\label{fig:fm:the-asset-manager-and-the-fund:stress}
\end{omfigure}

The pro rata slice keeps the fund's shape but costs the stayers 5.4\% on the hedge-fund ladder and 8.8\% on the illiquid one, because a slice of
every bucket must be sold within a week. \omterm{def:fm:the-asset-manager-and-the-fund:gate}{Swing pricing} moves the cost to the redeemers and leaves the liquidity problem where it is. A gate
buys time: it pays 10\% now, defers 20\%, and leaves the manager weeks to sell the slower assets at a normal cost. The illiquid fund has no good
option: whatever it does, its first 8\% of redemptions exhaust everything it can sell in a week. The largest redemption it can meet within seven
days at a cost below 1\% is 8\% of its NAV; the Form PF ladder's is 51.8\%.

\section{What investors ask for: side letters and managed accounts}

\begin{definition}[Side letter, most-favoured-nation clause]\label{def:fm:the-asset-manager-and-the-fund:side}
A \emph{side letter}\index{side letter} is an agreement between a fund (or its manager) and one investor that varies the fund's terms for that
investor: lower fees, more information, better liquidity, a capacity commitment. A \emph{most-favoured-nation clause}\index{most-favoured-nation
clause} is a promise that the investor will be offered any better terms given to another investor of the same or smaller size.
\end{definition}

A \omterm{def:fm:the-asset-manager-and-the-fund:side}{side letter} that gives one investor better liquidity creates a mismatch among investors: it redeems first and leaves the others with a less
liquid fund, which is why regulators and consultants ask whether any investor has preferential liquidity. Large investors that want full
transparency, their own liquidity and no exposure to other investors' flows use a \omterm{def:fm:the-asset-manager-and-the-fund:ucits}{separately managed account} instead: the manager trades the
investor's own account on the same model, and the investor's \omterm{def:fm:the-asset-manager-and-the-fund:admin}{custodian} holds the assets. The manager then runs several copies of one strategy, and
must allocate trades among them fairly, which its compliance function tests (chapter 16).

\begin{remark}[Liquidity is a property of the pair]\label{rem:fm:the-asset-manager-and-the-fund:pair}
A fund is not liquid or illiquid; its assets and its terms are liquid or not relative to each other. The same portfolio is sound in a vehicle with
quarterly dealing and 90 days' notice and fragile in one with daily dealing. A manager who changes the portfolio's liquidity, as the regulator
found in 2019, changes the fund's promise without changing its documents.
\end{remark}

\section{Tutorial: one fund, two investors and a bad week}\label{tut:fm:the-asset-manager-and-the-fund:tut}

\begin{tutorial}
\textbf{Goal.} Compute who pays which fees, and what a large redemption does to the investors who stay. \textbf{End state:} the stress table and
\cref{fig:fm:the-asset-manager-and-the-fund:stress}.
\begin{tutsteps}
\item \textbf{Fees by investor.} \texttt{fm\_fund.free\_ride()} runs \cref{ex:fm:the-asset-manager-and-the-fund:free} through
\texttt{firm.fundterms.pooled\_vs\_series}, which wraps One Quant Book 1's \texttt{firm.fees}.
\item \textbf{\omterm{def:fm:the-asset-manager-and-the-fund:class}{Crystallisation}.} \texttt{fm\_fund.crystallisation()} averages the fee take over 1\,000 paths for monthly, quarterly and annual
\omterm{def:fm:the-asset-manager-and-the-fund:class}{crystallisation}.
\item \textbf{Ladders.} \texttt{hedge\_fund\_ladder()} turns Form PF's cumulative portfolio liquidity into buckets; \texttt{illiquid\_ladder()}
builds the 8\%-in-seven-days fund.
\item \textbf{Stress.} \texttt{stress\_table()} meets a 30\% request under four policies (\cref{prop:fm:the-asset-manager-and-the-fund:liquidfirst}).
\end{tutsteps}
\end{tutorial}

\textbf{What to change next.} Run two redemption weeks in a row and watch the liquid-first fund's second week; give one investor a \omterm{def:fm:the-asset-manager-and-the-fund:side}{side letter} with
daily liquidity in a fund with monthly dealing.

\section{Build: the fund's terms}\label{bld:fm:the-asset-manager-and-the-fund:fundterms}

\begin{build}
\bfield{Purpose} The fund as its investors and administrator see it: fees by investor, \omterm{def:fm:the-asset-manager-and-the-fund:class}{crystallisation}, and what its liquidity terms do under
redemptions.
\bfield{Interface} \texttt{firm.fundterms}: \texttt{pooled\_vs\_series(returns, entries, terms)}; \texttt{crystallised(returns, perf, every)};
\texttt{Bucket(days, share, cost)}, \texttt{ladder\_share\_within}; \texttt{redeem(ladder, request, notice\_days, policy, gate, swing, fire)};
\texttt{max\_redemption(ladder, notice\_days, max\_cost)}.
\bfield{Rules} Fee arithmetic is One Quant Book 1's \texttt{firm.fees}, unchanged; without \omterm{def:fm:the-asset-manager-and-the-fund:gate}{swing pricing} the stayers bear the whole selling cost;
a gate defers, it never cancels.
\bfield{Acceptance tests} \texttt{code/firm/fundterms/tests/}: the free ride appears in the pooled fund and not in series; more frequent
\omterm{def:fm:the-asset-manager-and-the-fund:class}{crystallisation} never lowers the fee; liquid first is cheaper and leaves less liquidity than pro rata; the gate defers the rest; \omterm{def:fm:the-asset-manager-and-the-fund:gate}{swing pricing}
removes the stayers' dilution.
\bfield{Stretch} Multiple dealing dates with deferred requests queued; a \omterm{def:fm:the-asset-manager-and-the-fund:gate}{side pocket} that removes a bucket from the redeemable NAV; an investor
with a \omterm{def:fm:the-asset-manager-and-the-fund:side}{side letter}.
\end{build}

\begin{omsources}
\item Financial Conduct Authority, press releases of 11 April 2024 (Link Fund Solutions) and 5 August 2025 (Woodford Equity Income Fund).
\item US Securities and Exchange Commission, Division of Investment Management, \emph{Private Funds Statistics}, third calendar quarter 2025.
\item Directive 2009/65/EC (UCITS), article 84; 17 CFR 270.22e-4 and 270.22c-1 (US liquidity risk management and swing pricing).
\end{omsources}

\section{Exercises}

\begin{exercise}[$\star$]\label{exo:fm:the-asset-manager-and-the-fund:1}
In \cref{ex:fm:the-asset-manager-and-the-fund:free}, compute each investor's performance fee under the pooled share price and under series
accounting.
\end{exercise}

\begin{exercise}[$\star$]\label{exo:fm:the-asset-manager-and-the-fund:2}
From Form PF's 2025Q3 figures, what shares of NAV could qualifying hedge funds' investors withdraw, and their portfolios sell, within 30 days?
\end{exercise}

\begin{exercise}[$\star$]\label{exo:fm:the-asset-manager-and-the-fund:3}
Qualifying hedge funds reported \$121 billion gated, \$103 billion side-pocketed and \$25 billion suspended, on NAV of \$4\,750 billion. Express each
as a share of NAV.
\end{exercise}

\begin{exercise}[$\star\star$]\label{exo:fm:the-asset-manager-and-the-fund:4}
A fund can sell 20\% of its NAV within seven days and receives a 15\% redemption, met liquid first. What share of the remaining fund can be sold
within seven days?
\end{exercise}

\begin{exercise}[$\star\star$]\label{exo:fm:the-asset-manager-and-the-fund:5}
Why does a performance fee crystallised monthly take more than one crystallised annually on the same returns? Give a two-month example.
\end{exercise}

\begin{exercise}[$\star\star$]\label{exo:fm:the-asset-manager-and-the-fund:6}
In \cref{fig:fm:the-asset-manager-and-the-fund:stress}, why does the pro rata policy cost the stayers more than liquid first, yet leave them a more
liquid fund?
\end{exercise}

\begin{exercise}[$\star\star\star$]\label{exo:fm:the-asset-manager-and-the-fund:7}
\emph{Coding.} Run \texttt{firm.fundterms.redeem} on the illiquid ladder for a 5\% request, then a second 5\% request on the ladder left after the
first, both liquid first with seven days' notice. What share can be sold within seven days after each?
\end{exercise}

\begin{exercise}[$\star\star\star$]\label{exo:fm:the-asset-manager-and-the-fund:8}
\emph{Find the flaw.} ``Our fund can sell 60\% of its assets within a week and our investors have monthly liquidity with 30 days' notice, so we
have no liquidity risk.''
\end{exercise}

\section{Problem: The Friday Queue}

\begin{problem}[{Weekend problem --- the Friday queue}]\label{pb:fm:the-asset-manager-and-the-fund:1}
On Friday an open-ended fund with weekly dealing and seven days' notice receives redemption requests for 30\% of its NAV. You run its liquidity.

\medskip\noindent\textbf{Part I --- The vehicle.}
\begin{enumerate}
\item Define a \omterm{def:fm:the-asset-manager-and-the-fund:master}{master--feeder structure} and the roles of the general and \omterm{def:fm:the-asset-manager-and-the-fund:gp}{limited partners}.
\item What do the administrator and the \omterm{def:fm:the-asset-manager-and-the-fund:admin}{custodian} do, and why must they be independent of the manager?
\item Define a \omterm{def:fm:the-asset-manager-and-the-fund:notice}{redemption notice period}, a lock-up, a gate, a \omterm{def:fm:the-asset-manager-and-the-fund:gate}{side pocket} and \omterm{def:fm:the-asset-manager-and-the-fund:gate}{swing pricing}.
\item What does the \omterm{def:fm:the-asset-manager-and-the-fund:ucits}{UCITS} directive require about redemptions?
\end{enumerate}

\medskip\noindent\textbf{Part II --- The public record.}
\begin{enumerate}[resume]
\item Give Form PF's investor and portfolio liquidity at seven days for 2025Q3.
\item Is the hedge-fund industry mismatched in aggregate? At which horizon is the gap smallest in relative terms?
\item Give the shares of qualifying hedge-fund NAV gated, side-pocketed and suspended.
\item What did the regulator find about the fund suspended in 2019, and what is the status of its decisions?
\end{enumerate}

\medskip\noindent\textbf{Part III --- The queue.}
\begin{enumerate}[resume]
\item State and prove \cref{prop:fm:the-asset-manager-and-the-fund:liquidfirst}.
\item For the Form PF ladder, give the stayers' cost and the share sellable within seven days afterwards under liquid first and pro rata.
\item The same for the illiquid fund.
\item What does a 10\% gate pay, defer, and leave liquid, for each ladder?
\item What does \omterm{def:fm:the-asset-manager-and-the-fund:gate}{swing pricing} change, and what not?
\item Give the largest redemption each ladder can meet within seven days at a cost below 1\%.
\end{enumerate}

\medskip\noindent\textbf{Part IV --- The fees.}
\begin{enumerate}[resume]
\item Define \omterm{def:fm:the-asset-manager-and-the-fund:eq}{equalisation} and compute the free ride in \cref{ex:fm:the-asset-manager-and-the-fund:free}.
\item Give the mean fee take for monthly, quarterly and annual \omterm{def:fm:the-asset-manager-and-the-fund:class}{crystallisation}.
\item Why is a \omterm{def:fm:the-asset-manager-and-the-fund:side}{side letter} with better liquidity a liquidity risk for the other investors?
\item Define a \omterm{def:fm:the-asset-manager-and-the-fund:mismatch}{liquidity mismatch}.
\item State the \emph{named result}: the largest redemption each fund can meet within its notice period at a cost below 1\%, and the dilution the
stayers bear without \omterm{def:fm:the-asset-manager-and-the-fund:gate}{swing pricing} when 30\% is met liquid first.
\item In two sentences, what should the illiquid fund's board have done before Friday?
\end{enumerate}
\end{problem}

\section{Interview questions}

\begin{interviewq}[$\star$ \iqroles{researcher}]\label{iq:fm:the-asset-manager-and-the-fund:1}
What is a \omterm{def:fm:the-asset-manager-and-the-fund:master}{master--feeder structure} for, and who owns the portfolio?
\end{interviewq}

\begin{interviewq}[$\star$ \iqroles{trader, risk}]\label{iq:fm:the-asset-manager-and-the-fund:2}
What is a \omterm{def:fm:the-asset-manager-and-the-fund:gate}{redemption gate}, and why might investors prefer a fund that has one?
\end{interviewq}

\begin{interviewq}[$\star\star$ \iqroles{researcher}]\label{iq:fm:the-asset-manager-and-the-fund:3}
An investor buys into a pooled fund 10\% below its high-water mark. Explain the free ride and how a fund prevents it.
\end{interviewq}

\begin{interviewq}[$\star\star$ \iqroles{risk}]\label{iq:fm:the-asset-manager-and-the-fund:4}
You must meet a large redemption. Why is selling the most liquid assets first not neutral for the investors who stay?
\end{interviewq}

\begin{interviewq}[$\star\star$ \iqroles{trader}]\label{iq:fm:the-asset-manager-and-the-fund:5}
What does \omterm{def:fm:the-asset-manager-and-the-fund:gate}{swing pricing} do, and what problem does it not solve?
\end{interviewq}

\begin{interviewq}[$\star\star\star$ \iqroles{risk, researcher}]\label{iq:fm:the-asset-manager-and-the-fund:6}
How would you measure a fund's \omterm{def:fm:the-asset-manager-and-the-fund:mismatch}{liquidity mismatch}, and what data would you ask the manager for?
\end{interviewq}
